\documentclass[UTF8,12pt]{ctexart}
\usepackage[a4paper,margin=24mm]{geometry}
\usepackage{amsmath,amssymb,bm,graphicx,adjustbox}
\newcommand{\fitmath}[1]{\begin{adjustbox}{max width=\linewidth}$\displaystyle #1$\end{adjustbox}}
\setlength{\parindent}{0pt}
\setlength{\parskip}{.7em}
\sloppy
\allowdisplaybreaks
\begin{document}
\section*{2002 年考研数学三 · 参考答案}
\subsection*{2002年真题参考答案}

\subsection*{一、填空题}

（1）\(\displaystyle \dfrac{\displaystyle 1}{\displaystyle 1-2a}\)。（2）\(\displaystyle \int_0^{1/2}\mathrm dx\displaystyle \int_{x^2}^{x}f(x,\allowbreak y)\,\mathrm dy\)。（3）\(\displaystyle -1\)。（4）\(\displaystyle -0.02\)。（5）\(\displaystyle \bar X-1\)。


\subsection*{二、选择题}

（1）B。（2）无正确答案。（考题有误：幂级数 \(\displaystyle \sum_{n=1}^{\infty}a_nx^n\) 的收敛半径等于 \(\displaystyle R\)，推不出 \(\displaystyle \lim_{n\to\infty}|a_{n+1}/a_n|=1/R\)，因为 \(\displaystyle \lim_{n\to\infty}|a_{n+1}/a_n|\) 可能不存在。增加“设 \(\displaystyle \lim_{n\to\infty}|a_{n+1}/a_n|\) 和 \(\displaystyle \lim_{n\to\infty}|b_{n+1}/b_n|\) 均存在”的条件后，选项A正确。）（3）D。（4）B。（5）C。


\subsection*{三、}

\(\displaystyle \dfrac{\displaystyle \pi}{\displaystyle 6}\)。


\subsection*{四、}

\(\displaystyle \mathrm du=\left(f_x^{\prime}+f_z^{\prime}\dfrac{\displaystyle x+1}{\displaystyle z+1}e^{x-z}\right)\mathrm dx+\left(f_y^{\prime}-f_z^{\prime}\dfrac{\displaystyle y+1}{\displaystyle z+1}e^{y-z}\right)\mathrm dy\)。


\subsection*{五、}

\(\displaystyle -2\sqrt{1-x}\arcsin\sqrt x+2\sqrt x+C\)，其中 \(\displaystyle C\) 为任意常数。


\subsection*{六、}

（1）\(\displaystyle V_1=\dfrac{\displaystyle 4\pi}{\displaystyle 5}(32-a^5)\)，\(\displaystyle V_2=\pi a^4\)。（2）当 \(\displaystyle a=1\) 时，\(\displaystyle V_1+V_2\) 取得最大值，为 \(\displaystyle \dfrac{\displaystyle 129}{\displaystyle 5}\pi\)。


\subsection*{七、}

（1）证明略。（分别计算 \(\displaystyle y^{\prime}\)、\(\displaystyle y^{\prime\prime}\)，代入所需证明的等式。）（2）\(\displaystyle y(x)=\dfrac{\displaystyle 2}{\displaystyle 3}e^{-x/2}\cos\dfrac{\displaystyle \sqrt3}{\displaystyle 2}x+\dfrac{\displaystyle 1}{\displaystyle 3}e^x\;(-\infty<x<+\infty)\)。


\subsection*{八、}

证明略。（可利用连续函数的最值定理与介值定理。）


\subsection*{九、}

当 \(\displaystyle a\ne b\) 且 \(\displaystyle a\ne(1-n)b\) 时，方程组仅有零解。当 \(\displaystyle a=b\) 时，方程组有无穷多解，其解为 \(\displaystyle x=k_1\alpha_1+k_2\alpha_2+\cdots+k_{n-1}\alpha_{n-1}\)，其中 \(\displaystyle k_1,\allowbreak k_2,\allowbreak \ldots,\allowbreak k_{n-1}\) 为任意常数，\(\displaystyle \alpha_1=(-1,\allowbreak 1,\allowbreak 0,\allowbreak \ldots,\allowbreak 0)^{\mathrm T}\)，\(\displaystyle \alpha_2=(-1,\allowbreak 0,\allowbreak 1,\allowbreak \ldots,\allowbreak 0)^{\mathrm T}\)，\(\displaystyle \ldots\)，\(\displaystyle \alpha_{n-1}=(-1,\allowbreak 0,\allowbreak 0,\allowbreak \ldots,\allowbreak 1)^{\mathrm T}\)。当 \(\displaystyle a=(1-n)b\) 时，方程组有无穷多解，其解为 \(\displaystyle x=k(1,\allowbreak 1,\allowbreak \ldots,\allowbreak 1)^{\mathrm T}\)，其中 \(\displaystyle k\) 为任意常数。


\subsection*{十、}

（1）\(\displaystyle A\) 的全部特征值为 \(\displaystyle -2,\allowbreak -2,\allowbreak 0\)。（2）当 \(\displaystyle k>2\) 时，\(\displaystyle A+kE\) 为正定矩阵。


\subsection*{十一、}

（1）\(\displaystyle (X,\allowbreak Y)\) 的联合概率分布为 \(\displaystyle \begin{array}{c|cccc}(X,\allowbreak Y)&(-1,\allowbreak -1)&(-1,\allowbreak 1)&(1,\allowbreak -1)&(1,\allowbreak 1)\\\hline P&\dfrac{\displaystyle 1}{\displaystyle 4}&0&\dfrac{\displaystyle 1}{\displaystyle 2}&\dfrac{\displaystyle 1}{\displaystyle 4}\end{array}\)。（2）\(\displaystyle 2\)。


\subsection*{十二、}

\(\displaystyle F(y)=\begin{cases}0,\allowbreak &y<0,\allowbreak \\1-e^{-y/5},\allowbreak &0\le y<2,\allowbreak \\1,\allowbreak &y\ge2.\end{cases}\)

\end{document}
